% Pursuing Stacks (cotes 134-2 to 134-8): archive cote 134-3, site page 72 (PDF page 73), typescript page 72.
% typescript. Draft compared: ps2.tex:3717-3785.
% AI-transcribed by Claude Opus 5.5 (claude-opus-5-5, adaptive thinking, effort high), Message Batches API, batch ps1 (msgbatch_01Rv5hCaNTtUt4JHx1yqoHPL), 2026-11-12; not yet checked by eye.
\scanpage{134-3/72}
\mspage{72}

\margin{\unsure{$\ast$ cf th p. VIII}}ces in (Cat) is strongly saturated\hand{\unsure{(*)}}\hand{.} However, from the known relation
between Cat and an elementary modelizer $A\hat{}$, it follows that $W_{(\mathrm{Cat})}$ is
strongly saturated iff $W_A \hand{\subset} \mathrm{Fl}(A\hand{\hat{}})$ is. This implies that it is enough to
prove strong saturation in \uline{one} elementary modelizer $A\hat{}$, to deduce it in
all others, as well as in (Cat). However, in the case at least when
$A = \hand{\Delta}$, it is indeed "well-known" that \inserted{the} weak equivalences in $A\hat{}$ satisfy the
strong saturation condition. In\deleted{t} terms of Quillen's set-up, it follows
from the fact that $\hand{\Delta}\hat{}$ is a "closed model category" , and prop.1 in I 5.5.
of Quillen's expos\'e. The only thing which is not quite understood, not by
me at any rate, is why $\hand{\Delta}\hat{}$ \uline{is} indeed a closed model category - Quillen\hand{'}s
proof it seems relies strongly on typical simplicial techniques. I'll have
to look if the present set-up will suggest a more conceptual proof, valid
possibly for any test category (or at least, any strict test category).
I'll have to come back upon this later\hand{.} \deleted{xaftar} For the time being, I feel
a greater urge still to understand about the relationship between different
test categories\hand{,} - also\hand{,} I did not really finish with my review of what
may be viewed as \inserted{the} "standart\hand{"} test categories\deleted{"} such as $\hand{\Delta}$, $\hand{\square}$ and their vari-%
ants.

25.3.

42. In the notes last time I made clear what finally has turned out for
me to correspond to the appelation of a "modelizer", as prompted by the
internal logics of the situations I was looking at, in terms of the infor-%
mation available to me. I should by then have added what a \uline{model-preser-}%
\uline{ving} \inserted{(or \uline{modelizing})} \uline{functor} between modelizers (M,W) and (M',W') has turned out to mean,
which has undergone a corresponding change with respect \xout{xx} to what I first
contemplated calling by that name (\hand{\S.28}). Namely here it turned out that
I should be more stringent for this notion, replacing the condition
$f(W) \hand{\subset} W'$ by the stronger one
\[ W = f^{-1}(W') \quad , \]
and moreover, of course, still demanding that the corresponding functor
\[ W^{-1}M \hand{\to} W'^{-1}M' \]
should be an equivalence. It is by now established \deleted{that}, with \inserted{the} \xout{xxx} excep-%
tion of just \inserted{the two} \xout{xxx}, \inserted{first, that,} all the functors occuring in the diagram on p.31 are
indeed model\xout{xxxxx} preserving, namely it is so for $\hand{\alpha}$, $\hand{\beta}$, $\hand{\xi}$, $\hand{\eta}$ - and also
the right adjoints to $\hand{\xi}$, $\hand{\eta}$, as expected by then. It should be more or \xout{xx}
less trivial that the first of the functors in this diagram, namely
the canonical inclusion (Ord)$\hand{\longrightarrow}$(Preord), as well as the left adjoint
from (Preord) to (Ord), are model preserving - except for the fact that it
has not been yet established that these two categories are indeed modeli\xout{xx}
zers (for the natural notion of weak equivalences, induce\unsure{\over{s}{d}} from (Cat))\hand{,}
